\documentclass[11pt]{article}
\usepackage[margin=0.85in]{geometry}
\usepackage[T1]{fontenc}
\usepackage{lmodern,amsmath,amssymb,amsthm,mathtools,microtype}
\usepackage[hidelinks]{hyperref}
\usepackage{enumitem}
\setlist{nosep,leftmargin=*}
\setlength{\emergencystretch}{2em}
\newtheorem{theorem}{Theorem}
\newtheorem{lemma}[theorem]{Lemma}
\newtheorem{corollary}[theorem]{Corollary}
\newcommand{\End}{\operatorname{End}}
\newcommand{\Q}{\mathbb Q}
\newcommand{\R}{\mathbb R}
\newcommand{\C}{\mathbb C}
\newcommand{\N}{\mathsf N}
\title{\vspace{-1.5em}A degree-polynomial proof of spectral comparison\\
in every codimension on an abelian variety}
\author{\small Research note prepared for Christopher D. Long}
\date{\small 27 November 2026}
\begin{document}
\maketitle
\vspace{-1.2em}
\begin{abstract}
We retain the mixed-intersection terms in the shifted-degree method.
Real nonnegativity of the degree polynomial gives conjugate pairing;
testing at an exponential scale forces mixed-degree growth at
strict spectral gaps; log-concavity fills repeated-modulus blocks.
This gives a concise proof of Hu's established equality
\(\chi_{2k}=\lambda_k\) and his odd-degree formula, with a quantitative
byproduct weaker than recent general growth theorems. The proof uses neither the Albert
classification nor Deligne's weight theorem. It extends our earlier
first-degree argument. Its exposed-vertex strategy has close precedents
in Xie's comparison proof. No theorem or method priority is claimed.
\end{abstract}

\section{Statement and inputs}
Let \(A\) be an abelian variety of dimension \(g>0\) over an algebraically
closed field \(K\), and let \(f:A\to A\) be an isogeny. Fix a symmetric ample
class \(H\), a prime \(\ell\ne\operatorname{char}K\), and an abstract field embedding
\(\iota:\overline{\Q}_\ell\hookrightarrow\C\). All absolute values on
cohomology use this embedding; no continuity is intended. Write
\(\N^k(A)_\R\) for codimension-\(k\) cycles modulo numerical equivalence
with real coefficients. All divisor-class equalities are numerical. Set
\[
 \delta_k(n)=\frac{((f^n)^*H)^k\cdot H^{g-k}}{H^g},\quad
 \chi_i=\rho(f^*\mid H^i_{\mathrm{\acute et}}(A,\Q_\ell)),\quad
 \lambda_k=\rho(f^*\mid\N^k(A)_\R).
 \tag{1}
\]
\begin{theorem}[Known comparison theorem; proof below]\label{main}
For \(0\le k\le g\),
\[
 \lim_{n\to\infty}\delta_k(n)^{1/n}=\lambda_k=\chi_{2k}.
 \tag{2}
\]
For \(0\le k<g\),
\[
 \chi_{2k+1}=\sqrt{\lambda_k\lambda_{k+1}}.
 \tag{3}
\]
No semisimplicity assumption is needed.
\end{theorem}
The even-degree equality is Hu's theorem \cite[Theorem 1.9]{Hu}, and
the odd-degree formula is already \cite[Remark 5.5]{Hu}.
Xie proves the even comparison for surjective endomorphisms of all smooth projective varieties
\cite{Xie}, using Deligne's theorem. Our purpose is a short abelian-specific
argument with classical prerequisites.

We use the following classical inputs.
\begin{enumerate}
\item The characteristic polynomial on the Tate module is the degree
polynomial:
\[
 P_u(t)=\det(t-u\mid V_\ell A),\qquad
 P_u(a/b)=b^{-2g}\deg([a]-[b]u)\quad(a\in\mathbb Z,\ b>0).
 \tag{4}
\]
Here \(\deg u=0\) for a homomorphism that is not an isogeny;
the identity is valid even when the characteristic divides \(b\).
\item The theorem of the cube gives the numerical parallelogram identity;
pullbacks of \(H\) are nef; mixed intersections of nef classes are
nonnegative; and \((u^*H)^g=\deg(u)H^g\).
\item \(H^i(A)\simeq\bigwedge^i H^1(A)\), equivariantly, and cycle classes
are compatible with pullback and intersection. Tate twists are suppressed;
all actions are pullbacks by geometric endomorphisms.
\item For nef \(L,H\), the mixed intersections are log-concave:
\[
 (L^kH^{g-k})^2\ge
 (L^{k-1}H^{g-k+1})(L^{k+1}H^{g-k-1}).
 \tag{5}
\]
\end{enumerate}
Input (5) is the Khovanskii--Teissier inequality. For ample classes it
follows from the surface Hodge-index theorem on a general complete
intersection of \(k-1\) multiples of \(L\) and \(g-k-1\) multiples of \(H\);
nef classes follow by continuity. These inputs do not use the comparison
theorem or finite-field RH.

\section{Parity and the mixed-degree inequality}
\begin{lemma}[Degree positivity]\label{parity}
For \(u\in\End(A)\), \(P_u(t)\ge0\) for every real \(t\).
Every real root has even multiplicity, and the complex roots
can be arranged in conjugate pairs.
\end{lemma}
\begin{proof}
Equation (4) gives nonnegativity at every rational number. Continuity
gives nonnegativity on \(\R\), excluding an odd-multiplicity real
root. Nonreal roots pair under conjugation since \(P_u\in\Q[t]\).
\end{proof}
This proves the parity conclusion of \cite[Theorem 1.1]{Hu} directly
from the degree identity stated there as Theorem 2.3.
Write the root moduli of \(f\), in descending order, as
\[
 s_1,s_1,s_2,s_2,\ldots,s_g,s_g,\qquad s_1\ge\cdots\ge s_g>0,
 \quad U_k=\prod_{j=1}^k s_j^2,\quad U_0=1.
 \tag{6}
\]
Then \(\chi_{2k}=U_k\).

\begin{lemma}[Retaining all mixed degrees]\label{mixed}
For every real \(t\) and \(n\ge1\),
\[
 0\le P_{f^n}(t)\le
 2^g\sum_{j=0}^g\binom gj t^{2(g-j)}\delta_j(n).
 \tag{7}
\]
\end{lemma}
\begin{proof}
For rational \(t=a/b\), let
\[
 X=b^{-2}([a]-[b]f^n)^*H,\qquad
 Y=b^{-2}([a]+[b]f^n)^*H.
\]
Both are nef, and \(X+Y=2(t^2H+(f^n)^*H)\).
Consequently \(X^g\le(X+Y)^g\).
Divide by \(H^g\), apply (4), and expand.
Both sides of the scalar inequality are polynomials in \(t\),
so rational density proves the real case.
\end{proof}

\section{Recovering every exponential growth rate}
Each \(\delta_j(n)\) is a fixed matrix coefficient of \((f^*)^n\) on
\(H^{2j}(A)\), obtained by pairing \((f^*)^n[H^j]\) with
\([H^{g-j}]/H^g\). Finite-dimensional matrix growth gives
\[
 \delta_j(n)\le C_j n^{m_j-1}U_j^n,\qquad
 \limsup_{n\to\infty}\delta_j(n)^{1/n}\le U_j,
 \tag{8}
\]
where \(m_j\) is the largest Jordan-block size among eigenvalues of modulus
\(U_j\). Constants absorb the subdominant eigenvalues.

\paragraph{Strict gaps.}
Suppose \(s_k>s_{k+1}\), and choose \(s_{k+1}<R<s_k\).
Set \(t_n=2R^n\). Exactly \(2k\) roots have modulus greater than \(R\), and
\[
 \frac{P_{f^n}(2R^n)}
 {U_k^nR^{2n(g-k)}}\longrightarrow 2^{2(g-k)}.
 \tag{9}
\]
Indeed \(P_{f^n}(2R^n)\ge0\); in its product of absolute values, factors
above \(R\) are asymptotic to \(|\alpha|^n\), and factors below \(R\)
to \(2R^n\).
The sequence \(U_jR^{2(g-j)}\) has its unique maximum at \(j=k\),
since the ratio of successive terms is \(s_{j+1}^2/R^2\).
By (8), all terms \(j\ne k\) in (7) are exponentially smaller than (9).
Rearranging (7) gives
\[
 \delta_k(n)\ge c_k U_k^n\quad\text{for all sufficiently large }n.
 \tag{10}
\]
Thus \(\delta_k(n)^{1/n}\to U_k\) at every strict gap.
The endpoints are exact:
\(\delta_0(n)=1=U_0^n\) and
\(\delta_g(n)=(\deg f)^n=U_g^n\).

\paragraph{Repeated moduli.}
Let \(s_{a+1}=\cdots=s_b\) be a maximal constant block.
Its endpoint indices \(a,b\) are strict-gap indices or endpoints.
Repeated use of (5) gives, for \(a\le k\le b\),
\[
 \delta_k(n)\ge
 \delta_a(n)^{(b-k)/(b-a)}
 \delta_b(n)^{(k-a)/(b-a)}.
 \tag{11}
\]
Since \(U_k=U_a^{(b-k)/(b-a)}U_b^{(k-a)/(b-a)}\), the endpoint lower
bounds imply (10) throughout the block. Together with (8), this
proves existence of all limits in (2) and their equality to \(U_k\).
No existence theorem for dynamical-degree limits is being assumed.
The device of tilting until a vertex is the unique maximum, then filling
linear segments by concavity, also occurs in \cite[\S2.5]{Xie} and
\cite[Lemma 2.2 and the proof of Theorem 1.5]{HTX}. Here the input that enables it is (7), rather than
a Frobenius-weight comparison.

\paragraph{Numerical spectral radii.}
The same \(\delta_k(n)\) is a matrix coefficient on \(\N^k(A)_\R\),
so its established growth gives \(U_k\le\lambda_k\).
Conversely, \(\N^k(A)_\Q\otimes\Q_\ell\) is an equivariant quotient of
the subspace of \(H^{2k}(A,\Q_\ell)\) spanned by algebraic cycle classes:
a cohomologically zero linear combination has zero intersection with
every complementary cycle, hence is numerically zero.
Therefore \(\lambda_k\le\chi_{2k}=U_k\).
This does not assume that numerical and homological equivalence coincide.
Finally \(\chi_{2k+1}=U_ks_{k+1}=\sqrt{U_kU_{k+1}}\), proving (3).
\hfill\(\square\)

\section{Consequences and boundaries}
\paragraph{Quantitative consequence.}
There are constants \(c_k,C_k>0\) with
\[
 c_k\lambda_k^n\le\delta_k(n)
 \le C_k n^{m_k-1}\lambda_k^n\qquad(n\ge1).
 \tag{12}
\]
The finitely many initial \(n\) are absorbed into the constants.
These constants may depend on the fixed \(A,f,H,k\) and realization.
When the action on \(H^{2k}\) has no nontrivial Jordan block at its
spectral radius, the degree sequence is comparable to \(\lambda_k^n\).
No optimal universal bound for \(m_k\) is asserted. This is a byproduct,
not a new best growth estimate: Hu--Jiang \cite[Theorem 1.2,
Proposition 5.1]{HJ} prove the stronger relation
\(\deg_k(f^n)\asymp\lambda_k^n n^{\mu_k}\), with
\(\mu_k\le2k(g-k)\), for normal projective varieties in arbitrary
characteristic. Strict log-concavity forces \(\mu_k=0\) there.

\paragraph{Polarized and finite-field cases.}
If \(f^*H=dH\), then \(\delta_k(n)=d^{kn}\).
Equation (2), or its \(k=1\) case together with \(\deg f=d^g\),
gives \(|\alpha|=\sqrt d\) on \(H^1\).
Exterior powers give modulus \(d^{i/2}\) on \(H^i(A)\).
Frobenius over \(\mathbb F_q\) has \(d=q\) for an
\(\mathbb F_q\)-defined ample class. A curve's Jacobian gives
RH for curves and global function fields.
Frobenius here means pullback by the Frobenius endomorphism,
not the untwisted Galois action on divisor classes.

\paragraph{Translations and non-surjective maps.}
Translation acts trivially on cohomology and numerical classes, so
the theorem applies to all surjective self-morphisms.
For a group endomorphism that is not an isogeny, let \(2h\) be the number
of nonzero characteristic roots.
For \(k>h\), the operator on \(H^{2k}\) is nilpotent,
so \(\delta_k(n)=0\) eventually and both spectral radii are zero.
For the remaining terms, run the strict-gap argument on \(0,\ldots,h\),
exposing \(h\) by \(0<R<s_h\).
Log-concavity is used only inside the positive block.
The case \(h=0\) is immediate.

\begin{corollary}[Two maps with a common polarization]
If \(u,v\in\End(A)\), \(u^*H=aH\), \(v^*H=bH\), and \(a,b\ge0\), then
\[
 (\sqrt a-\sqrt b)^2H\preceq(u-v)^*H
 \preceq(\sqrt a+\sqrt b)^2H,
 \tag{13}
\]
and consequently
\[
 |\sqrt a-\sqrt b|^{2g}\le\deg(u-v)
 \le(\sqrt a+\sqrt b)^{2g}.
 \tag{14}
\]
Here \(D\preceq E\) means \(E-D\) is nef. The maps need not commute.
\end{corollary}
\begin{proof}
Write \((tu-v)^*H=(at^2+b)H-tB\).
For every curve \(C\), nefness for rational \(t\) gives a nonnegative
quadratic, so \(|B\cdot C|\le2\sqrt{ab}(H\cdot C)\).
This is (13). Monotonicity of top intersections gives (14).
\end{proof}

\paragraph{Relation to the recent short Weil-II proof.}
The Alp\"oge--Hexagon--Zhang argument uses generic affine weak Lefschetz,
finite-field Fourier estimates, and tensor amplification for arbitrary
lisse and constructible sheaves \cite{Hexagon,Zhang}.
Here the group law and degree polynomial restrict the scope to abelian
varieties. This proof strengthens our earlier first-degree argument to
all codimensions; it does not prove the arbitrary-sheaf weight theorem.
The established abelian theorem and odd-degree formula are Hu's;
Xie's theorem has broader scope. The useful presentation here is the
combination of (7) with the existing exposed-vertex/concavity strategy.
No historical shortest-proof or proof-novelty claim is made.

\begin{thebibliography}{9}\small
\bibitem{Hu}
F. Hu, \emph{Eigenvalues and dynamical degrees of self-maps on abelian
varieties}, J. Algebraic Geom. \textbf{33} (2024), 265--293.
arXiv:1909.12296v3, Theorems 1.1, 1.9, and 2.3.
\url{https://arxiv.org/abs/1909.12296}.
\bibitem{Milne}
J. S. Milne, \emph{The Riemann Hypothesis over Finite Fields:
From Weil to the Present Day}, 2015, \S1.
\url{https://www.jmilne.org/math/xnotes/pRH.pdf}.
\bibitem{Hexagon}
\emph{A short proof of Deligne's bound for Frobenius eigenvalues,
after Alp\"oge}, contributed by F. Alderliefste,
Hexagon 2611.00789v2, 24 November 2026.
\url{https://hexagonmath.org/pdf/2611.00789v2}.
\bibitem{Zhang}
D. Zhang, \emph{Notes on an elementary proof of Weil II from Hexagon},
arXiv:2611.30749, 2026.
\url{https://arxiv.org/abs/2611.30749}.
\bibitem{Xie}
J. Xie, \emph{Numerical spectrums control cohomological spectrums},
arXiv:2412.01216v2, 2025, Theorem 1.2 and \S2.5.
\url{https://arxiv.org/abs/2412.01216}.
\bibitem{HTX}
F. Hu, T. T. Truong, and J. Xie, \emph{Tate's question, Standard conjecture D,
semisimplicity and Dynamical degree comparison conjecture},
arXiv:2503.09432v1, 2025, Lemma 2.2.
\url{https://arxiv.org/abs/2503.09432}.
\bibitem{HJ}
F. Hu and C. Jiang, \emph{Uniform Tur\'an estimates and sharp bounds for
degree and mixed orbit growth}, arXiv:2609.29877v1,
24 September 2026.
\url{https://arxiv.org/abs/2609.29877}.
\end{thebibliography}
\end{document}
